% !TEX root = CoeneDylanBP.tex
%%=============================================================================
%% Systeemontwerp
%%=============================================================================

\chapter{Systeemontwerp}
\label{ch:systeemontwerp}

Dit hoofdstuk beschrijft de drie fundamentele bouwstenen waarop de meetpipeline steunt: de trainingsdata, de automatische detectie van de trampolinehoekpunten en de homografietransformatie die pixelco{\"o}rdinaten omzet naar metrische posities op het trampolinedoek. Samen vormen deze componenten het fundament waarop de landingsdetectie en tijdsmeting in Hoofdstuk~\ref{ch:meetpipeline} worden gebouwd.

%----------------------------------------------------------------------

\section{Dataset en Dataverzameling}
\label{sec:dataset}

Omdat de initi{\"e}le opzet met een eigen cameraopstelling werd verlaten, werd overgeschakeld naar een aanpak op basis van bestaande wedstrijdbeelden. De dataset van \textcite{Lin2025}, die videomateriaal van competitieve trampolinesprongen bevat, werd gebruikt als primaire databron voor het trainen en evalueren van de modellen in de pipeline.

Deze dataset bestaat uit beelden van de Olympische Spelen en andere internationale wedstrijden, waarbij de camera niet volledig gefixeerd was maar meebewoog met de actie, vergelijkbaar met hoe een smartphone-applicatie in de praktijk zou worden ingezet.

Een direct gevolg van het werken met pre-opgenomen beelden is dat er geen controle bestond over de camerahoeken, instellingen of de positie ten opzichte van de trampoline.

Uit de dataset werden gelabelde afbeeldingen geselecteerd voor twee afzonderlijke trainingsdoeleinden:
\begin{itemize}
    \item Het trainen van het hoekpuntdetectiemodel (Sectie~\ref{sec:trampolinedetectie}): afbeeldingen van trampolines waarbij de vier hoeken van het bed handmatig werden geannoteerd.
    \item Het trainen van de landingsclassifier (Sectie~\ref{subsec:landingsmoment-classifier}): uitgeknipte beelden van het trampolinebed, voorzien van het label \textit{Landed} of \textit{Not Landed}, afhankelijk van of een springer in contact was met het doek.
\end{itemize}

%----------------------------------------------------------------------

\section{Trampolinedetectie en Hoekpuntlocalisatie}
\label{sec:trampolinedetectie}

Voordat enige positiemeting mogelijk is, moeten de vier hoekpunten van het trampolinebed per frame gelokaliseerd worden. Deze hoekpunten vormen de invoer voor de homografietransformatie en bepalen daarmee de nauwkeurigheid van alle volgende stappen.

\subsection{Hoekpuntdetectie via YOLO-Pose}
\label{subsec:hoekpuntdetectie}

Voor de automatische localisatie van de hoekpunten werd een YOLO11n-pose model ingezet \autocite{Jocher_Ultralytics_YOLO_2023}. Dit model is toegelicht in Sectie~\ref{subsec:yolo} van de literatuurstudie. De \textit{pose}-variant van YOLO11 laat toe om naast een bounding box ook een vaste set keypoints te voorspellen. In dit geval werd het model geconfigureerd met vier keypoints die overeenkomen met de vier hoeken van het trampolinebed: linksboven, linksonder, rechtsboven en rechtsonder.

Het model werd getraind op een dataset van enkele honderden handmatig geannoteerde trampolineafbeeldingen, geselecteerd uit de dataset beschreven in Sectie~\ref{sec:dataset}. Dankzij de uiteenlopende camerahoeken en -afstanden in het wedstrijdmateriaal was het model na training in staat om de hoekpunten te detecteren ondanks de bewegende camera en wisselende perspectieven.

De nauwkeurigheid van het model was echter niet perfect. Bij frames waarbij het trampolinedoek sterk doorzakt, wat typisch het geval is op het moment van landing, vertoonden de gedetecteerde hoekpunten soms een zichtbare verschuiving ten opzichte van de werkelijke hoekposities. Omdat de volledige homografietransformatie op deze vier punten steunt, vertaalt elke fout in de hoekpuntdetectie zich rechtstreeks in een fout in de berekende landingspositie. Dit vormt de belangrijkste nauwkeurigheidsbeperking van de automatische aanpak.

\subsection{Alternatief: Manuele Annotatie}
\label{subsec:manuele-annotatie}

Als alternatief voor de automatische detectie werd een optie voorzien om de vier hoekpunten handmatig aan te duiden in het eerste frame van een videosegment. Dit sluit aan bij een scenario waarbij de cameraopstelling voor een bepaalde opname vast blijft en de hoekpunten slechts eenmalig bepaald hoeven te worden.

Manuele annotatie elimineert de detectiefout van het model volledig, waardoor de homografiematrix een stuk nauwkeuriger wordt. In de experimenten leidde dit tot een duidelijk betere landingspositiebepaling in vergelijking met de volledig geautomatiseerde aanpak. Het nadeel is uiteraard de schaalbaarheid: bij grote hoeveelheden video, of wanneer de camera doorheen de opname beweegt, is handmatig labelen niet praktisch haalbaar. Voor dit onderzoek fungeerde manuele annotatie als een bovengrens voor de haalbare nauwkeurigheid van de homografie.

%----------------------------------------------------------------------

\section{Homografietransformatie: van pixel naar meter}
\label{sec:homografie}

De pixelco{\"o}rdinaten $(u, v)$ uit het beeldvlak worden via een basistransformatie geprojecteerd naar fysieke co{\"o}rdinaten $(x, y)$ op het trampolinedoek, uitgedrukt in meters. Omdat een camera de trampoline meestal onder een hoek bekijkt, lijkt de rechthoekige mat op beeld vaak op een trapezium. Aangezien de vorm en afmetingen van een wedstrijdtrampoline gestandaardiseerd zijn (zie Sectie~\ref{subsec:trampolinetoestel}), kan er gebruik worden gemaakt van een \textit{planaire homografietransformatie} om dit perspectief te corrigeren.

Een homografie $H$ is een $3 \times 3$ matrix die de wiskundige relatie beschrijft tussen twee platte vlakken: het camerabeeld en het werkelijke trampolinedoek. Deze matrix dient als een brug die elke pixel in de video koppelt aan een specifieke plek op de mat.

De transformatie wordt bepaald aan de hand van vier corresponderende puntenparen:
\begin{enumerate}
    \item De vier hoekpunten van de trampolinemat worden in het videoframe (pixelco{\"o}rdinaten) bepaald via de methode beschreven in Sectie~\ref{sec:trampolinedetectie}.
    \item De re{\"e}le afmetingen van de trampoline (4,28\,m $\times$ 2,14\,m, zie Sectie~\ref{subsec:trampolinetoestel}) worden als doelco{\"o}rdinaten opgegeven.
\end{enumerate}

Met behulp van functies zoals \texttt{cv2.getPerspectiveTransform} of \texttt{cv2.findHomography} uit de OpenCV-bibliotheek wordt de matrix $H$ berekend. Zodra deze matrix bekend is, kan elke pixel $(u, v)$ die op het trampolinedoek valt, worden omgezet naar een exacte metrische co{\"o}rdinaat $(x, y)$ via de volgende berekening:

\[
    \begin{pmatrix} x' \\ y' \\ w \end{pmatrix} = H \cdot \begin{pmatrix} u \\ v \\ 1 \end{pmatrix}, \quad (x, y) = \left(\frac{x'}{w},\, \frac{y'}{w}\right)
\]

Hierbij zorgt de factor $w$ voor de schaling die nodig is om het perspectief-effect volledig te neutraliseren. De nauwkeurigheid van de homografie wordt tot slot ge{\"e}valueerd door visuele meetpunten of markeringen op bekende posities op de mat te projecteren. De berekende afstand wordt dan vergeleken met de werkelijk gemeten afstand. De verwachte maximale afwijking bedraagt slechts enkele centimeters. Deze kleine foutmarge is afhankelijk van de resolutie van de camera en hoe precies het model de hoekpunten van de mat kan bepalen.

% TODO: Figuur die de homografie toont + plot van de evaluatie
